% ============================================================
%  C++ Design Patterns — Building a Key-Value Store CLI
%  Compile: xelatex book.tex && xelatex book.tex
% ============================================================
\documentclass[11pt,a4paper]{book}

% --- Fonts & Korean ---
\usepackage{fontspec}
\usepackage{xeCJK}
\setCJKmainfont{AppleGothic}
\setmainfont{Palatino}
\setmonofont{Menlo}[Scale=0.85]

% --- Layout ---
\usepackage[margin=2.5cm]{geometry}
\usepackage{setspace}
\onehalfspacing

% --- Packages ---
\usepackage{listings}
\usepackage{xcolor}
\usepackage{tcolorbox}
\tcbuselibrary{skins,breakable}
\usepackage{enumitem}
\usepackage{hyperref}
\usepackage{titlesec}
\usepackage{fancyhdr}
\usepackage{graphicx}
\usepackage{booktabs}
\usepackage{longtable}
\usepackage{array}
\usepackage{tabularx}

% --- Colors ---
\definecolor{codebg}{HTML}{F5F5F0}
\definecolor{codeframe}{HTML}{CCCCCC}
\definecolor{keyword}{HTML}{0000FF}
\definecolor{string}{HTML}{A31515}
\definecolor{comment}{HTML}{008000}
\definecolor{type}{HTML}{267F99}
\definecolor{boxblue}{HTML}{1565C0}
\definecolor{boxgreen}{HTML}{2E7D32}
\definecolor{boxorange}{HTML}{E65100}
\definecolor{boxpurple}{HTML}{6A1B9A}
\definecolor{quizbg}{HTML}{FFF8E1}
\definecolor{summarybg}{HTML}{E3F2FD}

% --- Code listings ---
\lstset{
  basicstyle=\small\ttfamily,
  backgroundcolor=\color{codebg},
  frame=single,
  rulecolor=\color{codeframe},
  breaklines=true,
  breakatwhitespace=false,
  tabsize=4,
  showstringspaces=false,
  numbers=left,
  numberstyle=\tiny\color{gray},
  numbersep=8pt,
  aboveskip=10pt,
  belowskip=10pt,
  xleftmargin=15pt,
  framexleftmargin=15pt,
  columns=flexible,
  keepspaces=true,
  literate={→}{{\ensuremath{\rightarrow}}}1,
}

\lstdefinestyle{cpp}{
  language=C++,
  keywordstyle=\color{keyword}\bfseries,
  stringstyle=\color{string},
  commentstyle=\color{comment}\itshape,
  morekeywords={override,final,nullptr,constexpr,noexcept,auto,
                 decltype,static_assert,unique_ptr,shared_ptr,
                 make_unique,make_shared,optional,string_view},
  emph={std,map,string,vector,ifstream,ofstream,istringstream,
        cout,cerr,endl,cin,getline,optional,nullopt,unique_ptr},
  emphstyle=\color{type},
}

\lstdefinestyle{cmake}{
  language={},
  keywordstyle=\color{keyword}\bfseries,
  morekeywords={cmake_minimum_required,project,set,add_executable,
                target_compile_features,VERSION,LANGUAGES},
  stringstyle=\color{string},
  commentstyle=\color{comment}\itshape,
  morecomment=[l]{\#},
}

\lstdefinestyle{shell}{
  language=bash,
  keywordstyle=\color{keyword}\bfseries,
  stringstyle=\color{string},
  commentstyle=\color{comment}\itshape,
}

% --- tcolorbox styles ---
\newtcolorbox{summarybox}[1][]{
  colback=summarybg,
  colframe=boxblue,
  fonttitle=\bfseries,
  title={배운 것 정리 (What You Learned)},
  breakable,
  #1
}

\newtcolorbox{quizbox}[1][]{
  colback=quizbg,
  colframe=boxorange,
  fonttitle=\bfseries,
  title={퀴즈 (Quiz)},
  breakable,
  #1
}

\newtcolorbox{conceptbox}[1][]{
  colback=white,
  colframe=boxgreen,
  fonttitle=\bfseries,
  breakable,
  #1
}

\newtcolorbox{warningbox}[1][]{
  colback=white,
  colframe=boxorange,
  fonttitle=\bfseries,
  title={주의 (Warning)},
  breakable,
  #1
}

\newtcolorbox{pythonbox}[1][]{
  colback=white,
  colframe=boxpurple,
  fonttitle=\bfseries,
  title={Python과 비교},
  breakable,
  #1
}

% --- Header/Footer ---
\pagestyle{fancy}
\fancyhf{}
\fancyhead[LE]{\leftmark}
\fancyhead[RO]{\rightmark}
\fancyfoot[C]{\thepage}
\renewcommand{\headrulewidth}{0.4pt}

% --- Chapter title format ---
\titleformat{\chapter}[display]
  {\normalfont\huge\bfseries}
  {\chaptertitlename\ \thechapter}{20pt}{\Huge}

% --- Hyperref ---
\hypersetup{
  colorlinks=true,
  linkcolor=boxblue,
  urlcolor=boxblue,
  pdftitle={C++ Design Patterns: Building a Key-Value Store CLI},
  pdfauthor={},
}

% --- Convenience macros ---
\newcommand{\cpp}{\texttt{C++}}
\newcommand{\file}[1]{\texttt{#1}}
\newcommand{\code}[1]{\texttt{#1}}

% ============================================================
\begin{document}

% --- Title Page ---
\begin{titlepage}
\centering
\vspace*{3cm}
{\Huge\bfseries C++ Design Patterns\par}
\vspace{0.5cm}
{\LARGE Building a Key-Value Store CLI\par}
\vspace{1.5cm}
{\Large\itshape Python 개발자를 위한 Modern C++ 패턴 입문\par}
\vspace{3cm}
{\large mini-kvstore 프로젝트를 만들면서 배우는\\
RAII, Command, Factory, Strategy 패턴\par}
\vfill
{\large 2026\par}
\end{titlepage}

% --- Table of Contents ---
\tableofcontents

% ============================================================
\chapter{프로젝트 설정과 빌드 시스템}
\label{ch:setup}
% ============================================================

이 장에서는 C++ 프로젝트의 기본 구조를 만들고, CMake로 빌드하는 방법을 배웁니다.

\section{왜 빌드 시스템이 필요한가?}

Python에서는 \code{python main.py}로 바로 실행할 수 있습니다. C++에서는 \textbf{컴파일}이라는 단계가 필요합니다. 소스 코드(\file{.cpp})를 기계어로 바꾸는 과정입니다.

\begin{pythonbox}
\textbf{Python}: \code{python main.py} → 인터프리터가 바로 실행\\
\textbf{C++}: \code{g++ main.cpp -o main \&\& ./main} → 컴파일 후 실행\\[6pt]
파일이 여러 개면? → 빌드 시스템(CMake)이 컴파일 순서를 관리해줍니다.
\end{pythonbox}

\section{CMake 기초}

CMake는 C++ 빌드를 관리하는 도구입니다. \file{CMakeLists.txt}라는 설정 파일에 ``어떤 파일을 컴파일해서 어떤 프로그램을 만들지''를 적습니다.

\begin{conceptbox}[title={CMake의 핵심 세 줄}]
\begin{enumerate}
\item \code{cmake\_minimum\_required} — 최소 CMake 버전
\item \code{project} — 프로젝트 이름
\item \code{add\_executable} — 어떤 소스 파일로 어떤 실행 파일을 만들지
\end{enumerate}
\end{conceptbox}

\section{프로젝트 디렉토리 구조}

\begin{lstlisting}[style=shell,numbers=none]
mini-kvstore/
  CMakeLists.txt       # 빌드 설정
  src/
    main.cpp           # 프로그램 진입점
    store.h            # KVStore 선언 (이후 장에서)
    store.cpp           # KVStore 구현 (이후 장에서)
    command.h          # Command 인터페이스 (이후 장에서)
    command.cpp        # Command 구현 (이후 장에서)
    exporter.h         # Exporter 인터페이스 (이후 장에서)
    exporter.cpp       # Exporter 구현 (이후 장에서)
\end{lstlisting}

\section{Include Guard}

C++에서 헤더 파일(\file{.h})은 여러 곳에서 \code{\#include} 될 수 있습니다. 같은 파일이 두 번 포함되면 ``이미 정의됨'' 에러가 납니다. 이를 방지하는 것이 \textbf{include guard}입니다.

\begin{lstlisting}[style=cpp]
#ifndef STORE_H    // 아직 정의 안 됐으면
#define STORE_H    // 정의하고 아래 내용 포함

// ... 헤더 내용 ...

#endif             // 끝
\end{lstlisting}

\begin{pythonbox}
Python은 모듈 시스템이 자동으로 중복 import를 막아줍니다.\\
C++은 텍스트 기반 \code{\#include}이므로 직접 막아야 합니다.
\end{pythonbox}

\section{코드: CMakeLists.txt}

\begin{lstlisting}[style=cmake,title={\file{CMakeLists.txt}}]
cmake_minimum_required(VERSION 3.16)
project(mini-kvstore LANGUAGES CXX)

set(CMAKE_CXX_STANDARD 17)
set(CMAKE_CXX_STANDARD_REQUIRED ON)

add_executable(kvstore
    src/main.cpp
)
\end{lstlisting}

\textbf{설명:}
\begin{itemize}[nosep]
\item \code{VERSION 3.16} — 대부분의 시스템에 있는 안정 버전
\item \code{LANGUAGES CXX} — C++ 프로젝트임을 명시
\item \code{CMAKE\_CXX\_STANDARD 17} — C++17 사용 (optional, map, string\_view 등)
\item \code{add\_executable} — \file{src/main.cpp}를 컴파일해서 \code{kvstore} 실행 파일 생성
\end{itemize}

\section{코드: main.cpp (Hello World)}

\begin{lstlisting}[style=cpp,title={\file{src/main.cpp}}]
#include <iostream>

int main() {
    std::cout << "mini-kvstore v0.1" << std::endl;
    return 0;
}
\end{lstlisting}

\section{빌드 및 실행}

\begin{lstlisting}[style=shell,numbers=none]
mkdir build && cd build
cmake ..
make
./kvstore
# 출력: mini-kvstore v0.1
\end{lstlisting}

\begin{summarybox}
\begin{itemize}[nosep]
\item C++은 컴파일 → 링크 → 실행의 과정이 필요하다
\item CMake는 이 과정을 자동화하는 빌드 시스템이다
\item \code{cmake\_minimum\_required}, \code{project}, \code{add\_executable}이 핵심 세 명령어다
\item Include guard는 헤더 파일의 중복 포함을 방지한다
\item \code{std::cout}으로 표준 출력에 텍스트를 쓴다
\end{itemize}
\end{summarybox}

\begin{quizbox}
아래 CMakeLists.txt의 빈칸을 채우세요:

\begin{lstlisting}[style=cmake,numbers=none]
cmake_minimum_required(_______ 3.16)
_______(mini-kvstore LANGUAGES CXX)

set(CMAKE_CXX_STANDARD ___)
set(CMAKE_CXX_STANDARD_REQUIRED ON)

add_executable(kvstore
    _____________
)
\end{lstlisting}

\textbf{정답}: VERSION, project, 17, src/main.cpp
\end{quizbox}

% ============================================================
\chapter{RAII와 값 의미론 — KVStore 클래스}
\label{ch:raii}
% ============================================================

이 장에서는 C++의 가장 핵심적인 개념인 \textbf{RAII}와 \textbf{값 의미론}을 배우고, 이를 활용해 인메모리 Key-Value Store 클래스를 만듭니다.

\section{RAII (Resource Acquisition Is Initialization)}

RAII는 ``자원 획득은 초기화다''라는 뜻입니다. C++에서 가장 중요한 패턴 중 하나입니다.

\begin{conceptbox}[title={RAII 핵심 원리}]
\begin{enumerate}
\item 생성자에서 자원을 획득한다 (파일 열기, 메모리 할당 등)
\item 소멸자에서 자원을 해제한다 (파일 닫기, 메모리 반환 등)
\item 객체가 스코프를 벗어나면 \textbf{자동으로} 소멸자가 호출된다
\end{enumerate}
\end{conceptbox}

\begin{pythonbox}
Python은 가비지 컬렉터(GC)가 메모리를 관리합니다.\\
C++은 GC가 없습니다. 대신 RAII로 자원을 \textbf{결정론적으로} 관리합니다.\\[6pt]
Python의 \code{with open('f') as f:}도 RAII와 비슷한 개념입니다.\\
C++은 \textbf{모든 객체}가 기본적으로 이렇게 동작합니다.
\end{pythonbox}

\section{스택 vs 힙}

\begin{lstlisting}[style=cpp]
void example() {
    int x = 42;              // 스택: 함수 끝나면 자동 제거
    std::string s = "hello"; // 스택: 문자열 데이터는 내부적으로 힙에
}                             // x와 s 모두 여기서 자동 소멸
\end{lstlisting}

\begin{conceptbox}[title={스택 vs 힙}]
\textbf{스택(Stack)}: 함수 안에서 선언한 변수. 함수 끝나면 자동 제거. 빠르다.\\
\textbf{힙(Heap)}: \code{new}로 수동 할당. \code{delete}로 수동 해제. 느리다.\\[6pt]
원칙: 가능하면 스택을 쓴다. 힙이 필요하면 \code{unique\_ptr}을 쓴다 (5장에서).
\end{conceptbox}

\section{std::map — 연관 컨테이너}

\code{std::map<Key, Value>}는 키로 값을 찾는 자료구조입니다. Python의 \code{dict}와 같은 역할입니다.

\begin{lstlisting}[style=cpp]
#include <map>
#include <string>

std::map<std::string, std::string> data;
data["name"] = "Alice";        // 삽입
std::string v = data["name"];  // 조회: "Alice"
data.erase("name");            // 삭제
data.count("name");            // 존재 확인: 0 또는 1
data.size();                   // 개수
\end{lstlisting}

\section{const 참조 (const T\&)}

C++에서 함수 인자를 넘길 때, 기본적으로 \textbf{복사}가 일어납니다. 큰 객체를 복사하면 느립니다. \code{const T\&}는 ``복사 없이 읽기만 하겠다''는 뜻입니다.

\begin{lstlisting}[style=cpp]
// 나쁨: string이 복사됨
void bad(std::string key) { /* ... */ }

// 좋음: 복사 없이 원본을 읽기만 함
void good(const std::string& key) { /* ... */ }
\end{lstlisting}

\begin{pythonbox}
Python은 모든 것이 참조(reference)입니다. 복사가 기본이 아닙니다.\\
C++은 값(value)이 기본입니다. 참조를 원하면 \code{\&}를 명시해야 합니다.
\end{pythonbox}

\section{explicit 키워드}

\code{explicit}은 컴파일러의 \textbf{암묵적 타입 변환}을 막습니다. 생성자가 인자 1개일 때 반드시 붙여야 합니다.

\begin{lstlisting}[style=cpp]
class KVStore {
public:
    explicit KVStore(const std::string& filepath);
    // explicit 없으면: KVStore s = "data.txt"; 가 컴파일됨 (위험!)
    // explicit 있으면: KVStore s("data.txt"); 만 허용
};
\end{lstlisting}

\section{코드: store.h}

\begin{lstlisting}[style=cpp,title={\file{src/store.h}}]
#ifndef STORE_H
#define STORE_H

#include <map>
#include <string>
#include <vector>

class KVStore {
public:
    explicit KVStore(const std::string& filepath);

    void set(const std::string& key, const std::string& value);
    std::string get(const std::string& key) const;
    bool remove(const std::string& key);
    std::vector<std::pair<std::string, std::string>> list() const;
    std::size_t count() const;

private:
    std::string filepath_;
    std::map<std::string, std::string> data_;
};

#endif
\end{lstlisting}

\textbf{설명:}
\begin{itemize}[nosep]
\item \code{explicit KVStore(const std::string\& filepath)} — 파일 경로를 받는 생성자. 암묵적 변환 방지.
\item \code{const} 멤버 함수 (\code{get}, \code{list}, \code{count}) — 객체 상태를 변경하지 않음을 보장
\item \code{filepath\_} — 멤버 변수에 \code{\_} 접미사는 C++ 관례
\item \code{std::vector<std::pair<...>>} — key-value 쌍의 목록을 반환
\end{itemize}

\section{코드: store.cpp (인메모리 버전)}

\begin{lstlisting}[style=cpp,title={\file{src/store.cpp}}]
#include "store.h"
#include <stdexcept>

KVStore::KVStore(const std::string& filepath)
    : filepath_(filepath) {}

void KVStore::set(const std::string& key, const std::string& value) {
    data_[key] = value;
}

std::string KVStore::get(const std::string& key) const {
    auto it = data_.find(key);
    if (it == data_.end()) {
        throw std::runtime_error("Key not found: " + key);
    }
    return it->second;
}

bool KVStore::remove(const std::string& key) {
    return data_.erase(key) > 0;
}

std::vector<std::pair<std::string, std::string>> KVStore::list() const {
    return std::vector<std::pair<std::string, std::string>>(
        data_.begin(), data_.end());
}

std::size_t KVStore::count() const {
    return data_.size();
}
\end{lstlisting}

\textbf{설명:}
\begin{itemize}[nosep]
\item \code{: filepath\_(filepath)} — \textbf{멤버 초기화 리스트}. 생성자 본문 전에 멤버를 초기화.
\item \code{data\_.find(key)} — \code{map}에서 키를 검색. 없으면 \code{end()} 반환.
\item \code{it->second} — iterator가 가리키는 pair의 value (first는 key)
\item \code{data\_.erase(key)} — 삭제한 개수 반환 (0 또는 1)
\end{itemize}

\section{CMakeLists.txt 업데이트}

\begin{lstlisting}[style=cmake,title={\file{CMakeLists.txt} (업데이트)}]
cmake_minimum_required(VERSION 3.16)
project(mini-kvstore LANGUAGES CXX)

set(CMAKE_CXX_STANDARD 17)
set(CMAKE_CXX_STANDARD_REQUIRED ON)

add_executable(kvstore
    src/main.cpp
    src/store.cpp
)
\end{lstlisting}

\begin{summarybox}
\begin{itemize}[nosep]
\item RAII = 생성자에서 자원 획득, 소멸자에서 자원 해제, 스코프 벗어나면 자동 해제
\item 스택 변수는 함수 끝나면 자동 소멸. 가능하면 스택을 쓴다
\item \code{std::map}은 Python의 dict와 같다
\item \code{const T\&}는 복사 없이 읽기만 하겠다는 뜻
\item \code{explicit}은 암묵적 타입 변환을 막는다
\item 멤버 초기화 리스트(\code{: member\_(arg)})는 생성자에서 멤버를 초기화하는 방법
\end{itemize}
\end{summarybox}

\begin{quizbox}
아래 \file{store.h}의 빈칸을 채우세요:

\begin{lstlisting}[style=cpp,numbers=none]
#ifndef STORE_H
#define STORE_H

#include <___>     // map 헤더
#include <string>

class KVStore {
public:
    ________ KVStore(const std::string& filepath); // 암묵적 변환 방지

    void set(const std::string& key, const std::string& value);
    std::string get(const std::string& key) _____;  // 상태 불변 보장

private:
    std::string filepath_;
    std::___<std::string, std::string> data_;
};

#endif
\end{lstlisting}

\textbf{정답}: map, explicit, const, map
\end{quizbox}

% ============================================================
\chapter{파일 I/O와 std::optional — 영속성}
\label{ch:fileio}
% ============================================================

이 장에서는 파일 읽기/쓰기로 데이터를 저장하고, \code{std::optional}로 ``값이 없을 수 있음''을 표현합니다.

\section{파일 스트림도 RAII}

C++의 파일 스트림(\code{std::ifstream}, \code{std::ofstream})은 RAII 패턴을 따릅니다:

\begin{lstlisting}[style=cpp]
#include <fstream>

void write_example() {
    std::ofstream out("data.txt");  // 생성자: 파일 열기
    out << "hello" << std::endl;
}                                    // 소멸자: 파일 자동 닫기

void read_example() {
    std::ifstream in("data.txt");   // 생성자: 파일 열기
    std::string line;
    std::getline(in, line);         // 한 줄 읽기
}                                    // 소멸자: 파일 자동 닫기
\end{lstlisting}

\begin{pythonbox}
Python: \code{with open('f') as f:} → 블록 끝나면 닫힘\\
C++: 변수 스코프 끝나면 자동으로 닫힘 (with문이 필요 없음)
\end{pythonbox}

\section{std::optional}

\code{std::optional<T>}는 ``값이 있거나 없거나''를 표현합니다. Python의 \code{Optional[T]} 타입 힌트와 비슷하지만, 실제로 런타임에 값이 있는지 없는지를 확인합니다.

\begin{lstlisting}[style=cpp]
#include <optional>

std::optional<std::string> find(const std::string& key) {
    if (/* found */) return value;   // 값이 있을 때
    return std::nullopt;             // 값이 없을 때
}

// 사용법:
auto result = find("name");
if (result.has_value()) {
    std::cout << *result << std::endl;  // *로 값 꺼내기
}
// 또는:
if (result) {             // bool처럼 사용 가능
    std::cout << result.value() << std::endl;
}
\end{lstlisting}

\begin{warningbox}
\code{std::optional}에서 값이 없는데 \code{*result}나 \code{result.value()}를 호출하면 \textbf{undefined behavior} 또는 예외가 발생합니다. 반드시 먼저 확인하세요!
\end{warningbox}

\section{파일 형식}

간단한 텍스트 형식을 사용합니다. 한 줄에 하나의 key-value 쌍을 \code{TAB}으로 구분합니다:

\begin{lstlisting}[style=shell,numbers=none]
name	Alice
age	30
city	Seoul
\end{lstlisting}

TAB 구분을 선택한 이유: 키와 값에 공백이 들어갈 수 있지만, TAB은 덜 일반적이라 구분자로 적합합니다.

\section{코드: store.h (업데이트)}

\begin{lstlisting}[style=cpp,title={\file{src/store.h} (업데이트)}]
#ifndef STORE_H
#define STORE_H

#include <map>
#include <optional>
#include <string>
#include <vector>

class KVStore {
public:
    explicit KVStore(const std::string& filepath);

    void set(const std::string& key, const std::string& value);
    std::optional<std::string> get(const std::string& key) const;
    bool remove(const std::string& key);
    std::vector<std::pair<std::string, std::string>> list() const;
    std::size_t count() const;

    void save() const;
    void load();

private:
    std::string filepath_;
    std::map<std::string, std::string> data_;
};

#endif
\end{lstlisting}

\textbf{변경점:}
\begin{itemize}[nosep]
\item \code{get()} 반환 타입이 \code{std::string} → \code{std::optional<std::string>}로 변경
\item \code{save()}, \code{load()} 메서드 추가
\end{itemize}

\section{코드: store.cpp (파일 I/O 추가)}

\begin{lstlisting}[style=cpp,title={\file{src/store.cpp} (업데이트)}]
#include "store.h"
#include <fstream>
#include <sstream>
#include <stdexcept>

KVStore::KVStore(const std::string& filepath)
    : filepath_(filepath) {}

void KVStore::set(const std::string& key, const std::string& value) {
    data_[key] = value;
}

std::optional<std::string> KVStore::get(const std::string& key) const {
    auto it = data_.find(key);
    if (it == data_.end()) {
        return std::nullopt;
    }
    return it->second;
}

bool KVStore::remove(const std::string& key) {
    return data_.erase(key) > 0;
}

std::vector<std::pair<std::string, std::string>> KVStore::list() const {
    return std::vector<std::pair<std::string, std::string>>(
        data_.begin(), data_.end());
}

std::size_t KVStore::count() const {
    return data_.size();
}

void KVStore::save() const {
    std::ofstream out(filepath_);
    if (!out) {
        throw std::runtime_error("Cannot open file: " + filepath_);
    }
    for (const auto& [key, value] : data_) {
        out << key << '\t' << value << '\n';
    }
}

void KVStore::load() {
    std::ifstream in(filepath_);
    if (!in) {
        return;  // 파일 없으면 빈 상태로 시작
    }
    data_.clear();
    std::string line;
    while (std::getline(in, line)) {
        auto tab = line.find('\t');
        if (tab != std::string::npos) {
            std::string key = line.substr(0, tab);
            std::string value = line.substr(tab + 1);
            data_[key] = value;
        }
    }
}
\end{lstlisting}

\textbf{주요 패턴:}
\begin{itemize}[nosep]
\item \code{const auto\& [key, value]} — \textbf{구조적 바인딩}(C++17). pair를 분해하여 이름 부여.
\item \code{std::string::npos} — \code{find()}가 실패했을 때 반환하는 특수값
\item \code{line.substr(0, tab)} — 0부터 tab까지의 부분 문자열
\item \code{if (!out)} — 스트림의 bool 변환으로 파일 열기 실패 확인
\end{itemize}

\begin{summarybox}
\begin{itemize}[nosep]
\item \code{ifstream}/\code{ofstream}은 RAII — 스코프 벗어나면 자동으로 파일 닫힘
\item \code{std::optional<T>}는 ``값이 없을 수 있음''을 타입으로 표현
\item \code{std::nullopt}은 ``값 없음''을 나타내는 리터럴
\item 구조적 바인딩(\code{auto\& [k, v]})으로 pair를 분해할 수 있다
\item \code{string::find()}는 실패 시 \code{npos} 반환
\end{itemize}
\end{summarybox}

\begin{quizbox}
\code{save()} 함수의 빈칸을 채우세요:

\begin{lstlisting}[style=cpp,numbers=none]
void KVStore::save() const {
    std::_______ out(filepath_);   // 쓰기용 파일 스트림
    if (!out) {
        throw std::runtime_error("Cannot open file: " + filepath_);
    }
    for (const auto& [key, value] : _____) {
        out << key << '___' << value << '\n';  // 구분자
    }
}   // 여기서 파일이 자동으로 닫히는 이유: _____
\end{lstlisting}

\textbf{정답}: ofstream, data\_, \textbackslash t, RAII (소멸자가 호출되어 파일을 닫음)
\end{quizbox}

% ============================================================
\chapter{가상 인터페이스와 추상 클래스 — Command 인터페이스}
\label{ch:interface}
% ============================================================

이 장에서는 C++의 \textbf{가상 함수(virtual function)}와 \textbf{추상 클래스(abstract class)}를 배우고, Command 패턴의 기반이 되는 인터페이스를 만듭니다.

\section{다형성 (Polymorphism)}

\textbf{다형성}이란 같은 인터페이스로 다른 동작을 실행하는 것입니다.

\begin{pythonbox}
Python에서는 duck typing입니다:\\
\code{for cmd in commands: cmd.execute()} → 각 cmd가 다른 클래스여도 동작\\[6pt]
C++에서는 \textbf{명시적 인터페이스}가 필요합니다:\\
기반 클래스(base class)에 \code{virtual} 함수를 선언하고, 파생 클래스(derived class)에서 \code{override}합니다.
\end{pythonbox}

\section{순수 가상 함수 (Pure Virtual Function)}

\code{= 0}을 붙이면 ``구현이 없다''는 뜻입니다. 이런 함수가 하나라도 있으면 그 클래스는 \textbf{추상 클래스}가 됩니다. 추상 클래스의 인스턴스는 만들 수 없습니다.

\begin{lstlisting}[style=cpp]
class Shape {
public:
    virtual double area() const = 0;   // 순수 가상 함수
    virtual ~Shape() = default;        // 가상 소멸자
};

class Circle : public Shape {
public:
    explicit Circle(double r) : radius_(r) {}
    double area() const override {     // override로 명시
        return 3.14159 * radius_ * radius_;
    }
private:
    double radius_;
};
\end{lstlisting}

\section{가상 소멸자 (Virtual Destructor)}

기반 클래스 포인터로 파생 객체를 삭제할 때, 가상 소멸자가 없으면 파생 클래스의 소멸자가 호출되지 않습니다. 이것은 \textbf{메모리 누수}의 원인이 됩니다.

\begin{warningbox}
\textbf{규칙}: 가상 함수가 하나라도 있으면, 소멸자도 반드시 \code{virtual}로 선언하세요.\\
\code{virtual \~{}ClassName() = default;} — 대부분의 경우 이것으로 충분합니다.
\end{warningbox}

\section{override 키워드}

\code{override}는 ``이 함수가 기반 클래스의 가상 함수를 재정의한다''는 것을 컴파일러에게 알립니다. 기반 클래스에 해당 함수가 없으면 컴파일 에러가 납니다.

\begin{lstlisting}[style=cpp]
class Base {
public:
    virtual void execute() = 0;
};

class Derived : public Base {
public:
    void execute() override { /* ... */ }  // OK
    void exeucte() override { /* ... */ }  // 컴파일 에러! 오타 발견!
};
\end{lstlisting}

\section{코드: command.h (Set/Get 커맨드)}

\begin{lstlisting}[style=cpp,title={\file{src/command.h}}]
#ifndef COMMAND_H
#define COMMAND_H

#include "store.h"
#include <string>
#include <vector>

class Command {
public:
    virtual ~Command() = default;
    virtual void execute(KVStore& store) = 0;
};

class SetCommand : public Command {
public:
    SetCommand(const std::string& key, const std::string& value);
    void execute(KVStore& store) override;
private:
    std::string key_;
    std::string value_;
};

class GetCommand : public Command {
public:
    explicit GetCommand(const std::string& key);
    void execute(KVStore& store) override;
private:
    std::string key_;
};

#endif
\end{lstlisting}

\textbf{설명:}
\begin{itemize}[nosep]
\item \code{Command}는 추상 클래스 — \code{execute()}가 순수 가상 함수
\item \code{virtual \~{}Command() = default} — 가상 소멸자
\item \code{execute(KVStore\& store)} — KVStore 참조를 받아서 동작 수행
\item \code{SetCommand}, \code{GetCommand}가 \code{Command}를 상속하고 \code{override}
\end{itemize}

\section{코드: command.cpp (Set/Get 구현)}

\begin{lstlisting}[style=cpp,title={\file{src/command.cpp}}]
#include "command.h"
#include <iostream>

// --- SetCommand ---
SetCommand::SetCommand(const std::string& key, const std::string& value)
    : key_(key), value_(value) {}

void SetCommand::execute(KVStore& store) {
    store.set(key_, value_);
    store.save();
    std::cout << "OK" << std::endl;
}

// --- GetCommand ---
GetCommand::GetCommand(const std::string& key)
    : key_(key) {}

void GetCommand::execute(KVStore& store) {
    auto result = store.get(key_);
    if (result) {
        std::cout << *result << std::endl;
    } else {
        std::cerr << "Error: key not found: " << key_ << std::endl;
    }
}
\end{lstlisting}

\begin{summarybox}
\begin{itemize}[nosep]
\item \code{virtual ... = 0}은 순수 가상 함수 — 파생 클래스가 반드시 구현해야 함
\item 순수 가상 함수가 하나라도 있으면 추상 클래스 — 인스턴스 생성 불가
\item \code{override}는 재정의를 명시하여 오타를 컴파일 타임에 잡아줌
\item 가상 함수가 있으면 반드시 가상 소멸자를 선언
\item 인터페이스(추상 클래스)를 통해 다형성을 구현
\end{itemize}
\end{summarybox}

\begin{quizbox}
아래 Command 인터페이스의 빈칸을 채우세요:

\begin{lstlisting}[style=cpp,numbers=none]
class Command {
public:
    _______ ~Command() = default;       // 가상 소멸자
    _______ void execute(KVStore& store) ___; // 순수 가상 함수
};

class GetCommand : ______ Command {     // 상속
public:
    explicit GetCommand(const std::string& key);
    void execute(KVStore& store) ________; // 재정의 명시
};
\end{lstlisting}

\textbf{정답}: virtual, virtual, = 0, public, override
\end{quizbox}

% ============================================================
\chapter{Command 패턴 — 나머지 커맨드 구현}
\label{ch:command}
% ============================================================

이 장에서는 \textbf{Command 패턴}의 이론을 배우고, 나머지 커맨드(Delete, List, Count)를 구현합니다.

\section{Command 패턴이란?}

\begin{conceptbox}[title={Command 패턴}]
\textbf{핵심}: 요청(동작)을 \textbf{객체}로 캡슐화한다.\\[6pt]
왜?
\begin{itemize}[nosep]
\item 요청을 큐에 넣거나, 로그로 남기거나, 취소(undo)할 수 있다
\item 요청을 보내는 코드와 실행하는 코드를 분리한다
\item 새 명령을 추가할 때 기존 코드를 수정하지 않는다 (Open/Closed Principle)
\end{itemize}
\end{conceptbox}

\begin{pythonbox}
Python에서는 함수가 일급 객체이므로 함수 자체를 넘기면 됩니다:\\
\code{commands = \{"set": do\_set, "get": do\_get\}}\\[6pt]
C++에서도 함수 포인터나 lambda를 쓸 수 있지만, Command 패턴은 \textbf{상태(인자)를 함께 캡슐화}합니다. \code{SetCommand}는 key와 value를 멤버로 저장합니다.
\end{pythonbox}

\section{우리 프로젝트에서의 Command 패턴}

\begin{lstlisting}[style=shell,numbers=none]
./kvstore set name Alice   ->  SetCommand("name", "Alice").execute(store)
./kvstore get name         ->  GetCommand("name").execute(store)
./kvstore delete name      ->  DeleteCommand("name").execute(store)
./kvstore list             ->  ListCommand().execute(store)
./kvstore count            ->  CountCommand().execute(store)
\end{lstlisting}

CLI 인자를 파싱해서 적절한 Command 객체를 만들고, \code{execute()}를 호출합니다. 이 ``적절한 Command 객체를 만드는'' 부분은 다음 장(Factory 패턴)에서 다룹니다.

\section{코드: command.h (전체)}

\begin{lstlisting}[style=cpp,title={\file{src/command.h} (전체)}]
#ifndef COMMAND_H
#define COMMAND_H

#include "store.h"
#include <string>

class Command {
public:
    virtual ~Command() = default;
    virtual void execute(KVStore& store) = 0;
};

class SetCommand : public Command {
public:
    SetCommand(const std::string& key, const std::string& value);
    void execute(KVStore& store) override;
private:
    std::string key_;
    std::string value_;
};

class GetCommand : public Command {
public:
    explicit GetCommand(const std::string& key);
    void execute(KVStore& store) override;
private:
    std::string key_;
};

class DeleteCommand : public Command {
public:
    explicit DeleteCommand(const std::string& key);
    void execute(KVStore& store) override;
private:
    std::string key_;
};

class ListCommand : public Command {
public:
    void execute(KVStore& store) override;
};

class CountCommand : public Command {
public:
    void execute(KVStore& store) override;
};

#endif
\end{lstlisting}

\section{코드: command.cpp (전체)}

\begin{lstlisting}[style=cpp,title={\file{src/command.cpp} (전체)}]
#include "command.h"
#include <iostream>

// --- SetCommand ---
SetCommand::SetCommand(const std::string& key, const std::string& value)
    : key_(key), value_(value) {}

void SetCommand::execute(KVStore& store) {
    store.set(key_, value_);
    store.save();
    std::cout << "OK" << std::endl;
}

// --- GetCommand ---
GetCommand::GetCommand(const std::string& key)
    : key_(key) {}

void GetCommand::execute(KVStore& store) {
    auto result = store.get(key_);
    if (result) {
        std::cout << *result << std::endl;
    } else {
        std::cerr << "Error: key not found: " << key_ << std::endl;
    }
}

// --- DeleteCommand ---
DeleteCommand::DeleteCommand(const std::string& key)
    : key_(key) {}

void DeleteCommand::execute(KVStore& store) {
    if (store.remove(key_)) {
        store.save();
        std::cout << "OK" << std::endl;
    } else {
        std::cerr << "Error: key not found: " << key_ << std::endl;
    }
}

// --- ListCommand ---
void ListCommand::execute(KVStore& store) {
    auto entries = store.list();
    if (entries.empty()) {
        std::cout << "(empty)" << std::endl;
        return;
    }
    for (const auto& [key, value] : entries) {
        std::cout << key << "\t" << value << std::endl;
    }
}

// --- CountCommand ---
void CountCommand::execute(KVStore& store) {
    std::cout << store.count() << std::endl;
}
\end{lstlisting}

\section{CMakeLists.txt 업데이트}

\begin{lstlisting}[style=cmake,title={\file{CMakeLists.txt} (업데이트)}]
cmake_minimum_required(VERSION 3.16)
project(mini-kvstore LANGUAGES CXX)

set(CMAKE_CXX_STANDARD 17)
set(CMAKE_CXX_STANDARD_REQUIRED ON)

add_executable(kvstore
    src/main.cpp
    src/store.cpp
    src/command.cpp
)
\end{lstlisting}

\begin{summarybox}
\begin{itemize}[nosep]
\item Command 패턴 = 요청을 객체로 캡슐화
\item 각 Command 클래스는 실행에 필요한 인자를 멤버로 저장
\item \code{execute(KVStore\& store)}가 실제 동작을 수행
\item 새 명령 추가 시 기존 코드를 수정하지 않아도 된다 (OCP)
\item \code{ListCommand}, \code{CountCommand}는 인자가 필요 없으므로 멤버가 없다
\end{itemize}
\end{summarybox}

\begin{quizbox}
DeleteCommand를 직접 구현하세요:

\begin{lstlisting}[style=cpp,numbers=none]
class DeleteCommand : public Command {
public:
    ________ DeleteCommand(const std::string& key);
    void execute(KVStore& store) ________;
private:
    std::string ____;
};

// 구현:
DeleteCommand::DeleteCommand(const std::string& key)
    : ________(key) {}

void DeleteCommand::execute(KVStore& store) {
    if (store.______(key_)) {
        store.save();
        std::cout << "OK" << std::endl;
    } else {
        std::cerr << "Error: key not found: " << key_
                  << std::endl;
    }
}
\end{lstlisting}

\textbf{정답}: explicit, override, key\_, key\_(key), remove
\end{quizbox}

% ============================================================
\chapter{Factory 패턴과 unique\_ptr — CLI 파싱}
\label{ch:factory}
% ============================================================

이 장에서는 \textbf{Factory 패턴}과 \code{std::unique\_ptr}을 배우고, CLI 인자를 파싱하여 적절한 Command 객체를 생성하는 코드를 만듭니다.

\section{Factory 패턴이란?}

\begin{conceptbox}[title={Factory 패턴}]
\textbf{핵심}: 객체 생성 로직을 별도의 함수(또는 클래스)로 분리한다.\\[6pt]
왜?
\begin{itemize}[nosep]
\item 생성 로직이 복잡할 때 (조건에 따라 다른 타입 생성)
\item 호출자가 구체적인 타입을 몰라도 되게 할 때
\item 객체 생성 방식을 나중에 바꿀 수 있게 할 때
\end{itemize}
\end{conceptbox}

\section{std::unique\_ptr — 소유권}

\code{std::unique\_ptr<T>}는 ``이 포인터의 주인은 나 하나뿐''이라는 뜻입니다.

\begin{lstlisting}[style=cpp]
#include <memory>

// 생성: make_unique
auto ptr = std::make_unique<Circle>(5.0);

// 사용: 일반 포인터처럼
ptr->area();

// 소유권 이전: std::move
auto ptr2 = std::move(ptr);  // ptr은 이제 nullptr

// 소멸: 스코프 벗어나면 자동 delete
\end{lstlisting}

\begin{pythonbox}
Python에서는 모든 변수가 참조이고, GC가 메모리를 관리합니다.\\
C++에서 \code{unique\_ptr}은 GC 없이도 메모리 누수를 방지합니다.\\
소유권이 명확하므로 ``누가 이 객체를 지울 책임이 있는가?''가 타입으로 드러납니다.
\end{pythonbox}

\begin{warningbox}
\code{unique\_ptr}은 복사할 수 없습니다. 소유권은 하나뿐이니까요.\\
소유권을 넘기려면 반드시 \code{std::move()}를 사용해야 합니다.
\end{warningbox}

\section{std::move — 소유권 이전}

\code{std::move(x)}는 ``x의 자원을 가져가도 좋다''는 의미입니다. 이동 후 원래 변수는 비어있는(valid but unspecified) 상태가 됩니다.

\begin{lstlisting}[style=cpp]
auto cmd = std::make_unique<SetCommand>("k", "v");

// 함수에 소유권 넘기기:
void run(std::unique_ptr<Command> cmd) {
    cmd->execute(store);
}  // cmd가 스코프를 벗어나면서 자동 delete

run(std::move(cmd));  // 소유권 이전
// cmd는 이제 nullptr
\end{lstlisting}

\section{parseCommand 팩토리 함수}

우리 프로젝트에서 Factory 패턴의 구현입니다:

\begin{lstlisting}[style=shell,numbers=none]
./kvstore set name Alice
          ^^^ ^^^^ ^^^^^
          |   |    |
          |   |    args[1]  (argv[3])
          |   args[0]       (argv[2])
          command name      (argv[1])
\end{lstlisting}

\section{코드: main.cpp (완성)}

\begin{lstlisting}[style=cpp,title={\file{src/main.cpp} (완성)}]
#include "command.h"
#include "store.h"
#include <iostream>
#include <memory>
#include <string>

static void printUsage() {
    std::cerr << "Usage: kvstore <command> [args...]" << std::endl;
    std::cerr << "Commands:" << std::endl;
    std::cerr << "  set <key> <value>" << std::endl;
    std::cerr << "  get <key>" << std::endl;
    std::cerr << "  delete <key>" << std::endl;
    std::cerr << "  list" << std::endl;
    std::cerr << "  count" << std::endl;
}

static std::unique_ptr<Command> parseCommand(int argc, char* argv[]) {
    if (argc < 2) {
        return nullptr;
    }

    std::string cmd = argv[1];

    if (cmd == "set" && argc == 4) {
        return std::make_unique<SetCommand>(argv[2], argv[3]);
    }
    if (cmd == "get" && argc == 3) {
        return std::make_unique<GetCommand>(argv[2]);
    }
    if (cmd == "delete" && argc == 3) {
        return std::make_unique<DeleteCommand>(argv[2]);
    }
    if (cmd == "list" && argc == 2) {
        return std::make_unique<ListCommand>();
    }
    if (cmd == "count" && argc == 2) {
        return std::make_unique<CountCommand>();
    }

    return nullptr;
}

int main(int argc, char* argv[]) {
    auto cmd = parseCommand(argc, argv);
    if (!cmd) {
        printUsage();
        return 1;
    }

    KVStore store("kvstore.dat");
    store.load();
    cmd->execute(store);
    return 0;
}
\end{lstlisting}

\textbf{핵심 포인트:}
\begin{itemize}[nosep]
\item \code{parseCommand}는 \textbf{팩토리 함수} — CLI 인자를 보고 적절한 Command 객체를 생성
\item 반환 타입은 \code{unique\_ptr<Command>} — 기반 클래스 포인터로 다형성 활용
\item \code{std::make\_unique<SetCommand>(...)} — 힙에 SetCommand를 생성하고 unique\_ptr로 감싸기
\item \code{main()}은 어떤 Command가 생성되는지 몰라도 됨 — \code{cmd->execute()}만 호출
\item \code{static} 함수 — 이 파일에서만 사용 가능 (내부 링키지)
\end{itemize}

\begin{summarybox}
\begin{itemize}[nosep]
\item Factory 패턴 = 객체 생성 로직을 분리하여 호출자가 구체 타입을 몰라도 되게 함
\item \code{unique\_ptr<T>}는 소유권이 하나뿐인 스마트 포인터 — 스코프 벗어나면 자동 delete
\item \code{std::make\_unique<T>(...)}로 생성, \code{std::move()}로 소유권 이전
\item \code{unique\_ptr}은 복사 불가, 이동만 가능
\item Factory 함수가 \code{unique\_ptr<Base>}를 반환 → 다형성 + 자동 메모리 관리
\end{itemize}
\end{summarybox}

\begin{quizbox}
parseCommand 함수의 빈칸을 채우세요:

\begin{lstlisting}[style=cpp,numbers=none]
static std::__________<Command> parseCommand(
        int argc, char* argv[]) {
    if (argc < 2) return _______;

    std::string cmd = argv[1];

    if (cmd == "set" && argc == 4) {
        return std::____________<SetCommand>(
            argv[2], argv[3]);
    }
    if (cmd == "get" && argc == 3) {
        return std::____________<GetCommand>(argv[2]);
    }
    return nullptr;
}

int main(int argc, char* argv[]) {
    auto cmd = parseCommand(argc, argv);
    if (!cmd) { return 1; }

    KVStore store("kvstore.dat");
    store.load();
    cmd->_______(store);  // 다형성: 실제 타입의 execute 호출
    return 0;
}
\end{lstlisting}

\textbf{정답}: unique\_ptr, nullptr, make\_unique, make\_unique, execute
\end{quizbox}

% ============================================================
\chapter{Strategy 패턴 — Exporter 시스템}
\label{ch:strategy}
% ============================================================

이 장에서는 \textbf{Strategy 패턴}을 배우고, 데이터를 CSV나 JSON 형식으로 내보내는 Exporter 시스템을 만듭니다.

\section{Strategy 패턴이란?}

\begin{conceptbox}[title={Strategy 패턴}]
\textbf{핵심}: 알고리즘(전략)을 인터페이스 뒤에 숨기고, 실행 시점에 교체할 수 있게 한다.\\[6pt]
Command vs Strategy:
\begin{itemize}[nosep]
\item Command = ``무엇을 할 것인가''를 캡슐화 (set, get, delete...)
\item Strategy = ``어떻게 할 것인가''를 캡슐화 (CSV로? JSON으로?)
\end{itemize}
\end{conceptbox}

\begin{pythonbox}
Python 비유: 함수를 인자로 넘기는 것과 비슷합니다.\\
\code{def export\_data(data, formatter):}\\
\code{~~~~return formatter(data)}\\[6pt]
C++에서는 이것을 인터페이스(추상 클래스)로 구현합니다.
\end{pythonbox}

\section{std::ostream\& 추상화}

\code{std::ostream\&}는 ``출력 스트림에 대한 참조''입니다. \code{std::cout}(화면), \code{std::ofstream}(파일), \code{std::ostringstream}(문자열) 모두 \code{ostream}의 파생 클래스입니다.

\begin{lstlisting}[style=cpp]
// 이 함수는 화면, 파일, 문자열 어디에든 출력할 수 있다
void write(std::ostream& out) {
    out << "hello" << std::endl;
}

// 사용:
write(std::cout);              // 화면에
std::ofstream f("out.txt");
write(f);                      // 파일에
\end{lstlisting}

\section{코드: exporter.h}

\begin{lstlisting}[style=cpp,title={\file{src/exporter.h}}]
#ifndef EXPORTER_H
#define EXPORTER_H

#include <memory>
#include <ostream>
#include <string>
#include <utility>
#include <vector>

class Exporter {
public:
    virtual ~Exporter() = default;
    virtual void exportData(
        std::ostream& out,
        const std::vector<std::pair<std::string, std::string>>& data
    ) const = 0;
};

class CsvExporter : public Exporter {
public:
    void exportData(
        std::ostream& out,
        const std::vector<std::pair<std::string, std::string>>& data
    ) const override;
};

class JsonExporter : public Exporter {
public:
    void exportData(
        std::ostream& out,
        const std::vector<std::pair<std::string, std::string>>& data
    ) const override;
};

std::unique_ptr<Exporter> createExporter(const std::string& format);

#endif
\end{lstlisting}

\textbf{설명:}
\begin{itemize}[nosep]
\item \code{Exporter}는 Strategy 인터페이스 — \code{exportData()}가 순수 가상 함수
\item \code{CsvExporter}, \code{JsonExporter}가 구체적인 전략(Concrete Strategy)
\item \code{createExporter()}는 Factory 함수 — format 문자열로 적절한 Exporter 생성
\item Strategy + Factory 조합: 흔하게 사용되는 패턴
\end{itemize}

\section{코드: exporter.cpp}

\begin{lstlisting}[style=cpp,title={\file{src/exporter.cpp}}]
#include "exporter.h"
#include <stdexcept>

void CsvExporter::exportData(
    std::ostream& out,
    const std::vector<std::pair<std::string, std::string>>& data
) const {
    out << "key,value" << '\n';
    for (const auto& [key, value] : data) {
        out << key << ',' << value << '\n';
    }
}

void JsonExporter::exportData(
    std::ostream& out,
    const std::vector<std::pair<std::string, std::string>>& data
) const {
    out << "{\n";
    for (std::size_t i = 0; i < data.size(); ++i) {
        out << "  \"" << data[i].first
            << "\": \"" << data[i].second << "\"";
        if (i + 1 < data.size()) {
            out << ',';
        }
        out << '\n';
    }
    out << "}\n";
}

std::unique_ptr<Exporter> createExporter(const std::string& format) {
    if (format == "csv") {
        return std::make_unique<CsvExporter>();
    }
    if (format == "json") {
        return std::make_unique<JsonExporter>();
    }
    return nullptr;
}
\end{lstlisting}

\section{ExportCommand 추가}

command.h에 ExportCommand를 추가합니다:

\begin{lstlisting}[style=cpp,title={\file{src/command.h}에 추가}]
// (기존 #include 뒤에 추가)
#include "exporter.h"
#include <memory>

// ... (기존 Command 클래스들) ...

class ExportCommand : public Command {
public:
    explicit ExportCommand(std::unique_ptr<Exporter> exporter);
    void execute(KVStore& store) override;
private:
    std::unique_ptr<Exporter> exporter_;
};
\end{lstlisting}

\begin{lstlisting}[style=cpp,title={\file{src/command.cpp}에 추가}]
// --- ExportCommand ---
ExportCommand::ExportCommand(std::unique_ptr<Exporter> exporter)
    : exporter_(std::move(exporter)) {}

void ExportCommand::execute(KVStore& store) {
    auto data = store.list();
    exporter_->exportData(std::cout, data);
}
\end{lstlisting}

\textbf{주목}: \code{std::move(exporter)}로 소유권을 이전합니다. \code{unique\_ptr}은 복사할 수 없으므로 반드시 \code{move}해야 합니다.

\section{main.cpp에 export 명령 추가}

parseCommand 함수에 추가합니다:

\begin{lstlisting}[style=cpp,title={parseCommand에 추가}]
    if (cmd == "export" && argc == 3) {
        auto exporter = createExporter(argv[2]);
        if (!exporter) {
            std::cerr << "Error: unknown format: " << argv[2]
                      << std::endl;
            return nullptr;
        }
        return std::make_unique<ExportCommand>(std::move(exporter));
    }
\end{lstlisting}

사용법: \code{./kvstore export csv} 또는 \code{./kvstore export json}

\begin{summarybox}
\begin{itemize}[nosep]
\item Strategy 패턴 = 알고리즘을 인터페이스 뒤에 숨기고 교체 가능하게 함
\item \code{std::ostream\&}로 출력 대상을 추상화 (화면, 파일, 문자열)
\item Strategy + Factory 조합은 실무에서 자주 사용
\item \code{unique\_ptr}을 생성자로 넘길 때 \code{std::move()} 필수
\item 새로운 export 형식 추가 시 기존 코드 수정 불필요 (OCP)
\end{itemize}
\end{summarybox}

\begin{quizbox}
Strategy 패턴을 구현하세요:

\begin{lstlisting}[style=cpp,numbers=none]
class Exporter {
public:
    _______ ~Exporter() = default;
    virtual void exportData(
        std::_______ & out,   // 출력 스트림 타입
        const std::vector<...>& data
    ) const ___;              // 순수 가상
};

// Factory + Strategy 조합:
ExportCommand::ExportCommand(
        std::unique_ptr<Exporter> exporter)
    : exporter_(std::____(exporter)) {} // 소유권 이전
\end{lstlisting}

\textbf{정답}: virtual, ostream, = 0, move
\end{quizbox}

% ============================================================
\chapter{에러 처리와 방어적 프로그래밍 — 마무리}
\label{ch:error}
% ============================================================

이 장에서는 프로그램을 견고하게 만드는 에러 처리와 방어적 프로그래밍 기법을 적용합니다.

\section{에러 처리 전략}

우리 프로그램에서 발생할 수 있는 에러:

\begin{itemize}[nosep]
\item 잘못된 인자 개수 (\code{./kvstore set} — value 없음)
\item 존재하지 않는 키 (\code{./kvstore get nonexistent})
\item 알 수 없는 명령어 (\code{./kvstore foo})
\item 파일 쓰기 실패 (디스크 꽉 참, 권한 없음)
\item 빈 키 (\code{./kvstore set "" value})
\item 알 수 없는 export 형식 (\code{./kvstore export xml})
\end{itemize}

\section{입력 검증}

신뢰할 수 없는 입력은 시스템 경계(system boundary)에서 검증합니다. 우리 프로그램에서 시스템 경계는 \code{main()} — CLI 인자를 받는 지점입니다.

\begin{lstlisting}[style=cpp]
// 빈 키 검증을 parseCommand에 추가
if (cmd == "set" && argc == 4) {
    std::string key = argv[2];
    if (key.empty()) {
        std::cerr << "Error: key cannot be empty" << std::endl;
        return nullptr;
    }
    return std::make_unique<SetCommand>(key, argv[3]);
}
\end{lstlisting}

\section{const 정확성 (const correctness)}

\begin{conceptbox}[title={const correctness 규칙}]
\begin{enumerate}[nosep]
\item 변경하지 않는 매개변수: \code{const T\&}
\item 변경하지 않는 멤버 함수: \code{... const}
\item 변경하지 않는 지역 변수: \code{const auto}
\end{enumerate}
``가능한 한 많은 것을 const로'' — 의도를 명확히 하고 실수를 방지합니다.
\end{conceptbox}

\section{최종 코드: command.h}

\begin{lstlisting}[style=cpp,title={\file{src/command.h} (최종)}]
#ifndef COMMAND_H
#define COMMAND_H

#include "exporter.h"
#include "store.h"
#include <memory>
#include <string>

class Command {
public:
    virtual ~Command() = default;
    virtual void execute(KVStore& store) = 0;
};

class SetCommand : public Command {
public:
    SetCommand(const std::string& key, const std::string& value);
    void execute(KVStore& store) override;
private:
    std::string key_;
    std::string value_;
};

class GetCommand : public Command {
public:
    explicit GetCommand(const std::string& key);
    void execute(KVStore& store) override;
private:
    std::string key_;
};

class DeleteCommand : public Command {
public:
    explicit DeleteCommand(const std::string& key);
    void execute(KVStore& store) override;
private:
    std::string key_;
};

class ListCommand : public Command {
public:
    void execute(KVStore& store) override;
};

class CountCommand : public Command {
public:
    void execute(KVStore& store) override;
};

class ExportCommand : public Command {
public:
    explicit ExportCommand(std::unique_ptr<Exporter> exporter);
    void execute(KVStore& store) override;
private:
    std::unique_ptr<Exporter> exporter_;
};

#endif
\end{lstlisting}

\section{최종 코드: command.cpp}

\begin{lstlisting}[style=cpp,title={\file{src/command.cpp} (최종)}]
#include "command.h"
#include <iostream>

// --- SetCommand ---
SetCommand::SetCommand(const std::string& key, const std::string& value)
    : key_(key), value_(value) {}

void SetCommand::execute(KVStore& store) {
    store.set(key_, value_);
    store.save();
    std::cout << "OK" << std::endl;
}

// --- GetCommand ---
GetCommand::GetCommand(const std::string& key)
    : key_(key) {}

void GetCommand::execute(KVStore& store) {
    const auto result = store.get(key_);
    if (result) {
        std::cout << *result << std::endl;
    } else {
        std::cerr << "Error: key not found: " << key_ << std::endl;
    }
}

// --- DeleteCommand ---
DeleteCommand::DeleteCommand(const std::string& key)
    : key_(key) {}

void DeleteCommand::execute(KVStore& store) {
    if (store.remove(key_)) {
        store.save();
        std::cout << "OK" << std::endl;
    } else {
        std::cerr << "Error: key not found: " << key_ << std::endl;
    }
}

// --- ListCommand ---
void ListCommand::execute(KVStore& store) {
    const auto entries = store.list();
    if (entries.empty()) {
        std::cout << "(empty)" << std::endl;
        return;
    }
    for (const auto& [key, value] : entries) {
        std::cout << key << "\t" << value << std::endl;
    }
}

// --- CountCommand ---
void CountCommand::execute(KVStore& store) {
    std::cout << store.count() << std::endl;
}

// --- ExportCommand ---
ExportCommand::ExportCommand(std::unique_ptr<Exporter> exporter)
    : exporter_(std::move(exporter)) {}

void ExportCommand::execute(KVStore& store) {
    const auto data = store.list();
    exporter_->exportData(std::cout, data);
}
\end{lstlisting}

\section{최종 코드: main.cpp}

\begin{lstlisting}[style=cpp,title={\file{src/main.cpp} (최종)}]
#include "command.h"
#include "store.h"
#include <iostream>
#include <memory>
#include <string>

static void printUsage() {
    std::cerr << "Usage: kvstore <command> [args...]" << std::endl;
    std::cerr << "Commands:" << std::endl;
    std::cerr << "  set <key> <value>   Set a key-value pair" << std::endl;
    std::cerr << "  get <key>           Get value by key" << std::endl;
    std::cerr << "  delete <key>        Delete a key" << std::endl;
    std::cerr << "  list                List all pairs" << std::endl;
    std::cerr << "  count               Count entries" << std::endl;
    std::cerr << "  export <csv|json>   Export data" << std::endl;
}

static std::unique_ptr<Command> parseCommand(int argc, char* argv[]) {
    if (argc < 2) {
        return nullptr;
    }

    const std::string cmd = argv[1];

    if (cmd == "set") {
        if (argc != 4) {
            std::cerr << "Error: set requires <key> <value>"
                      << std::endl;
            return nullptr;
        }
        const std::string key = argv[2];
        if (key.empty()) {
            std::cerr << "Error: key cannot be empty" << std::endl;
            return nullptr;
        }
        return std::make_unique<SetCommand>(key, argv[3]);
    }

    if (cmd == "get") {
        if (argc != 3) {
            std::cerr << "Error: get requires <key>" << std::endl;
            return nullptr;
        }
        return std::make_unique<GetCommand>(argv[2]);
    }

    if (cmd == "delete") {
        if (argc != 3) {
            std::cerr << "Error: delete requires <key>" << std::endl;
            return nullptr;
        }
        return std::make_unique<DeleteCommand>(argv[2]);
    }

    if (cmd == "list") {
        if (argc != 2) {
            std::cerr << "Error: list takes no arguments" << std::endl;
            return nullptr;
        }
        return std::make_unique<ListCommand>();
    }

    if (cmd == "count") {
        if (argc != 2) {
            std::cerr << "Error: count takes no arguments" << std::endl;
            return nullptr;
        }
        return std::make_unique<CountCommand>();
    }

    if (cmd == "export") {
        if (argc != 3) {
            std::cerr << "Error: export requires <csv|json>"
                      << std::endl;
            return nullptr;
        }
        auto exporter = createExporter(argv[2]);
        if (!exporter) {
            std::cerr << "Error: unknown format: " << argv[2]
                      << " (use csv or json)" << std::endl;
            return nullptr;
        }
        return std::make_unique<ExportCommand>(std::move(exporter));
    }

    std::cerr << "Error: unknown command: " << cmd << std::endl;
    return nullptr;
}

int main(int argc, char* argv[]) {
    auto cmd = parseCommand(argc, argv);
    if (!cmd) {
        printUsage();
        return 1;
    }

    KVStore store("kvstore.dat");
    store.load();
    cmd->execute(store);
    return 0;
}
\end{lstlisting}

\section{최종 CMakeLists.txt}

\begin{lstlisting}[style=cmake,title={\file{CMakeLists.txt} (최종)}]
cmake_minimum_required(VERSION 3.16)
project(mini-kvstore LANGUAGES CXX)

set(CMAKE_CXX_STANDARD 17)
set(CMAKE_CXX_STANDARD_REQUIRED ON)

add_executable(kvstore
    src/main.cpp
    src/store.cpp
    src/command.cpp
    src/exporter.cpp
)
\end{lstlisting}

\begin{summarybox}
\begin{itemize}[nosep]
\item 입력 검증은 시스템 경계(main, CLI 인자)에서 한다
\item 내부 코드는 이미 검증된 데이터를 신뢰한다
\item 구체적인 에러 메시지를 stderr로 출력한다
\item const correctness: 변경하지 않는 것은 모두 const로
\item 각 명령의 인자 개수를 검증하고 구체적인 에러 메시지 출력
\end{itemize}
\end{summarybox}

\begin{quizbox}
방어적 프로그래밍 원칙을 적용하세요. 아래 코드에서 누락된 검증을 추가하세요:

\begin{lstlisting}[style=cpp,numbers=none]
if (cmd == "delete") {
    // Q1: argc가 정확히 몇이어야 하는가?
    if (argc != ___) {
        std::cerr << "Error: delete requires <key>"
                  << std::endl;
        return nullptr;
    }
    return std::make_unique<DeleteCommand>(argv[2]);
}

// Q2: 모든 if문에 해당하지 않은 경우, 무엇을 출력해야 하는가?
std::cerr << "Error: _______ command: " << cmd
          << std::endl;
return _______;
\end{lstlisting}

\textbf{정답}: 3, unknown, nullptr
\end{quizbox}

% ============================================================
\chapter{테스트와 검증 — test.sh}
\label{ch:test}
% ============================================================

이 장에서는 셸 스크립트로 프로그램의 통합 테스트(integration test)를 작성합니다.

\section{왜 셸 스크립트로 테스트하는가?}

CLI 프로그램은 ``입력 → 출력''이 명확합니다. 셸 스크립트에서 프로그램을 실행하고, 출력이 예상과 같은지 비교하면 됩니다.

\begin{pythonbox}
Python에서는 \code{pytest}를 쓰지만, CLI 도구 테스트에는 셸 스크립트가 더 직접적입니다.\\
우리가 테스트하는 것은 ``이 바이너리에 이 인자를 주면 이 출력이 나오는가?''입니다.
\end{pythonbox}

\section{테스트 패턴}

\begin{lstlisting}[style=shell,numbers=none]
# 기본 패턴:
output=$(./kvstore get name)
expected="Alice"
if [ "$output" != "$expected" ]; then
    echo "FAIL: expected '$expected', got '$output'"
    exit 1
fi
\end{lstlisting}

\section{코드: test.sh}

\begin{lstlisting}[style=shell,title={\file{test.sh}}]
#!/bin/bash
set -e

PASS=0
FAIL=0
BIN=./build/kvstore
DATA=kvstore.dat

pass() {
    PASS=$((PASS + 1))
    echo "  PASS: $1"
}

fail() {
    FAIL=$((FAIL + 1))
    echo "  FAIL: $1 - expected '$2', got '$3'"
}

check() {
    local desc="$1" expected="$2" actual="$3"
    if [ "$actual" = "$expected" ]; then
        pass "$desc"
    else
        fail "$desc" "$expected" "$actual"
    fi
}

check_stderr() {
    local desc="$1" pattern="$2"
    local output
    output=$($BIN $3 2>&1 || true)
    if echo "$output" | grep -q "$pattern"; then
        pass "$desc"
    else
        fail "$desc" "contains '$pattern'" "$output"
    fi
}

cleanup() {
    rm -f "$DATA"
}

echo "=== Building ==="
mkdir -p build && cd build && cmake .. -DCMAKE_BUILD_TYPE=Debug \
    > /dev/null 2>&1 && make > /dev/null 2>&1 && cd ..
echo "Build OK"
echo ""

# --- Test: Basic set/get ---
echo "=== Basic Operations ==="
cleanup
check "set key" "OK" "$($BIN set name Alice)"
check "get key" "Alice" "$($BIN get name)"
check "count after set" "1" "$($BIN count)"

# --- Test: Overwrite ---
echo "=== Overwrite ==="
$BIN set name Bob > /dev/null
check "get overwritten" "Bob" "$($BIN get name)"
check "count unchanged" "1" "$($BIN count)"

# --- Test: Multiple keys ---
echo "=== Multiple Keys ==="
$BIN set age 30 > /dev/null
$BIN set city Seoul > /dev/null
check "count multiple" "3" "$($BIN count)"

# --- Test: List ---
echo "=== List ==="
output=$($BIN list)
check "list contains age" "yes" \
    "$(echo "$output" | grep -q 'age' && echo yes || echo no)"
check "list contains name" "yes" \
    "$(echo "$output" | grep -q 'name' && echo yes || echo no)"

# --- Test: Delete ---
echo "=== Delete ==="
check "delete key" "OK" "$($BIN delete age)"
check "count after delete" "2" "$($BIN count)"

# --- Test: Get missing key ---
echo "=== Error Cases ==="
check_stderr "get missing key" "not found" "get nonexistent"

# --- Test: Delete missing key ---
check_stderr "delete missing key" "not found" "delete nonexistent"

# --- Test: No arguments ---
check_stderr "no arguments" "Usage" ""

# --- Test: Unknown command ---
check_stderr "unknown command" "unknown command" "foo"

# --- Test: Wrong arg count ---
check_stderr "set missing value" "requires" "set onlykey"

# --- Test: Persistence ---
echo "=== Persistence ==="
cleanup
$BIN set hello world > /dev/null
check "persist after restart" "world" "$($BIN get hello)"

# --- Test: Export CSV ---
echo "=== Export ==="
output=$($BIN export csv)
check "csv header" "yes" \
    "$(echo "$output" | grep -q 'key,value' && echo yes || echo no)"
check "csv data" "yes" \
    "$(echo "$output" | grep -q 'hello,world' && echo yes || echo no)"

# --- Test: Export JSON ---
output=$($BIN export json)
check "json format" "yes" \
    "$(echo "$output" | grep -q '"hello"' && echo yes || echo no)"

# --- Test: Unknown export format ---
check_stderr "unknown format" "unknown format" "export xml"

# --- Cleanup & Summary ---
cleanup
echo ""
echo "=== Results ==="
echo "Passed: $PASS"
echo "Failed: $FAIL"

if [ "$FAIL" -gt 0 ]; then
    exit 1
fi
echo "All tests passed!"
\end{lstlisting}

\textbf{핵심 포인트:}
\begin{itemize}[nosep]
\item \code{set -e} — 에러 발생 시 스크립트 즉시 중단
\item \code{check} 함수 — 테스트 패턴을 함수로 추출 (반복 제거)
\item \code{cleanup} — 각 테스트 섹션 전에 데이터 파일 삭제 (격리)
\item \code{\$(cmd)} — 명령의 stdout 캡처
\item \code{2>\&1} — stderr를 stdout으로 리다이렉트하여 에러 메시지도 캡처
\item \code{|| true} — 에러 종료 코드를 무시 (set -e에서 스크립트 중단 방지)
\end{itemize}

\begin{summarybox}
\begin{itemize}[nosep]
\item CLI 프로그램은 셸 스크립트로 통합 테스트를 작성하기 좋다
\item 테스트 패턴: 실행 → 출력 캡처 → 예상값과 비교
\item \code{set -e}로 실패 시 즉시 중단
\item helper 함수로 테스트 코드의 반복을 줄인다
\item 에러 케이스도 반드시 테스트한다
\end{itemize}
\end{summarybox}

\begin{quizbox}
셸 스크립트 테스트의 빈칸을 채우세요:

\begin{lstlisting}[style=shell,numbers=none]
#!/bin/bash
_____ -e    # 에러 시 즉시 중단

# stdout 캡처
output=____($BIN get name)

# 예상값 비교
if [ "$output" ____ "$expected" ]; then
    echo "PASS"
fi

# stderr도 캡처하려면:
output=$($BIN bad_cmd _____  || true)
\end{lstlisting}

\textbf{정답}: set, \$, =, 2>\&1
\end{quizbox}

% ============================================================
\appendix
% ============================================================

% ============================================================
\chapter{전체 코드 (Complete Code Reference)}
\label{app:code}
% ============================================================

이 부록에는 완성된 mini-kvstore의 전체 소스 코드가 있습니다.

\section{프로젝트 구조}

\begin{lstlisting}[style=shell,numbers=none]
mini-kvstore/
  CMakeLists.txt
  test.sh
  src/
    main.cpp
    store.h
    store.cpp
    command.h
    command.cpp
    exporter.h
    exporter.cpp
\end{lstlisting}

\section{CMakeLists.txt}

\begin{lstlisting}[style=cmake,title={\file{CMakeLists.txt}}]
cmake_minimum_required(VERSION 3.16)
project(mini-kvstore LANGUAGES CXX)

set(CMAKE_CXX_STANDARD 17)
set(CMAKE_CXX_STANDARD_REQUIRED ON)

add_executable(kvstore
    src/main.cpp
    src/store.cpp
    src/command.cpp
    src/exporter.cpp
)
\end{lstlisting}

\section{src/store.h}

\begin{lstlisting}[style=cpp,title={\file{src/store.h}}]
#ifndef STORE_H
#define STORE_H

#include <map>
#include <optional>
#include <string>
#include <vector>

class KVStore {
public:
    explicit KVStore(const std::string& filepath);

    void set(const std::string& key, const std::string& value);
    std::optional<std::string> get(const std::string& key) const;
    bool remove(const std::string& key);
    std::vector<std::pair<std::string, std::string>> list() const;
    std::size_t count() const;

    void save() const;
    void load();

private:
    std::string filepath_;
    std::map<std::string, std::string> data_;
};

#endif
\end{lstlisting}

\section{src/store.cpp}

\begin{lstlisting}[style=cpp,title={\file{src/store.cpp}}]
#include "store.h"
#include <fstream>
#include <sstream>
#include <stdexcept>

KVStore::KVStore(const std::string& filepath)
    : filepath_(filepath) {}

void KVStore::set(const std::string& key, const std::string& value) {
    data_[key] = value;
}

std::optional<std::string> KVStore::get(const std::string& key) const {
    auto it = data_.find(key);
    if (it == data_.end()) {
        return std::nullopt;
    }
    return it->second;
}

bool KVStore::remove(const std::string& key) {
    return data_.erase(key) > 0;
}

std::vector<std::pair<std::string, std::string>> KVStore::list() const {
    return std::vector<std::pair<std::string, std::string>>(
        data_.begin(), data_.end());
}

std::size_t KVStore::count() const {
    return data_.size();
}

void KVStore::save() const {
    std::ofstream out(filepath_);
    if (!out) {
        throw std::runtime_error("Cannot open file: " + filepath_);
    }
    for (const auto& [key, value] : data_) {
        out << key << '\t' << value << '\n';
    }
}

void KVStore::load() {
    std::ifstream in(filepath_);
    if (!in) {
        return;
    }
    data_.clear();
    std::string line;
    while (std::getline(in, line)) {
        auto tab = line.find('\t');
        if (tab != std::string::npos) {
            std::string key = line.substr(0, tab);
            std::string value = line.substr(tab + 1);
            data_[key] = value;
        }
    }
}
\end{lstlisting}

\section{src/command.h}

\begin{lstlisting}[style=cpp,title={\file{src/command.h}}]
#ifndef COMMAND_H
#define COMMAND_H

#include "exporter.h"
#include "store.h"
#include <memory>
#include <string>

class Command {
public:
    virtual ~Command() = default;
    virtual void execute(KVStore& store) = 0;
};

class SetCommand : public Command {
public:
    SetCommand(const std::string& key, const std::string& value);
    void execute(KVStore& store) override;
private:
    std::string key_;
    std::string value_;
};

class GetCommand : public Command {
public:
    explicit GetCommand(const std::string& key);
    void execute(KVStore& store) override;
private:
    std::string key_;
};

class DeleteCommand : public Command {
public:
    explicit DeleteCommand(const std::string& key);
    void execute(KVStore& store) override;
private:
    std::string key_;
};

class ListCommand : public Command {
public:
    void execute(KVStore& store) override;
};

class CountCommand : public Command {
public:
    void execute(KVStore& store) override;
};

class ExportCommand : public Command {
public:
    explicit ExportCommand(std::unique_ptr<Exporter> exporter);
    void execute(KVStore& store) override;
private:
    std::unique_ptr<Exporter> exporter_;
};

#endif
\end{lstlisting}

\section{src/command.cpp}

\begin{lstlisting}[style=cpp,title={\file{src/command.cpp}}]
#include "command.h"
#include <iostream>

SetCommand::SetCommand(const std::string& key, const std::string& value)
    : key_(key), value_(value) {}

void SetCommand::execute(KVStore& store) {
    store.set(key_, value_);
    store.save();
    std::cout << "OK" << std::endl;
}

GetCommand::GetCommand(const std::string& key)
    : key_(key) {}

void GetCommand::execute(KVStore& store) {
    const auto result = store.get(key_);
    if (result) {
        std::cout << *result << std::endl;
    } else {
        std::cerr << "Error: key not found: " << key_ << std::endl;
    }
}

DeleteCommand::DeleteCommand(const std::string& key)
    : key_(key) {}

void DeleteCommand::execute(KVStore& store) {
    if (store.remove(key_)) {
        store.save();
        std::cout << "OK" << std::endl;
    } else {
        std::cerr << "Error: key not found: " << key_ << std::endl;
    }
}

void ListCommand::execute(KVStore& store) {
    const auto entries = store.list();
    if (entries.empty()) {
        std::cout << "(empty)" << std::endl;
        return;
    }
    for (const auto& [key, value] : entries) {
        std::cout << key << "\t" << value << std::endl;
    }
}

void CountCommand::execute(KVStore& store) {
    std::cout << store.count() << std::endl;
}

ExportCommand::ExportCommand(std::unique_ptr<Exporter> exporter)
    : exporter_(std::move(exporter)) {}

void ExportCommand::execute(KVStore& store) {
    const auto data = store.list();
    exporter_->exportData(std::cout, data);
}
\end{lstlisting}

\section{src/exporter.h}

\begin{lstlisting}[style=cpp,title={\file{src/exporter.h}}]
#ifndef EXPORTER_H
#define EXPORTER_H

#include <memory>
#include <ostream>
#include <string>
#include <utility>
#include <vector>

class Exporter {
public:
    virtual ~Exporter() = default;
    virtual void exportData(
        std::ostream& out,
        const std::vector<std::pair<std::string, std::string>>& data
    ) const = 0;
};

class CsvExporter : public Exporter {
public:
    void exportData(
        std::ostream& out,
        const std::vector<std::pair<std::string, std::string>>& data
    ) const override;
};

class JsonExporter : public Exporter {
public:
    void exportData(
        std::ostream& out,
        const std::vector<std::pair<std::string, std::string>>& data
    ) const override;
};

std::unique_ptr<Exporter> createExporter(const std::string& format);

#endif
\end{lstlisting}

\section{src/exporter.cpp}

\begin{lstlisting}[style=cpp,title={\file{src/exporter.cpp}}]
#include "exporter.h"
#include <stdexcept>

void CsvExporter::exportData(
    std::ostream& out,
    const std::vector<std::pair<std::string, std::string>>& data
) const {
    out << "key,value" << '\n';
    for (const auto& [key, value] : data) {
        out << key << ',' << value << '\n';
    }
}

void JsonExporter::exportData(
    std::ostream& out,
    const std::vector<std::pair<std::string, std::string>>& data
) const {
    out << "{\n";
    for (std::size_t i = 0; i < data.size(); ++i) {
        out << "  \"" << data[i].first
            << "\": \"" << data[i].second << "\"";
        if (i + 1 < data.size()) {
            out << ',';
        }
        out << '\n';
    }
    out << "}\n";
}

std::unique_ptr<Exporter> createExporter(const std::string& format) {
    if (format == "csv") {
        return std::make_unique<CsvExporter>();
    }
    if (format == "json") {
        return std::make_unique<JsonExporter>();
    }
    return nullptr;
}
\end{lstlisting}

\section{src/main.cpp}

\begin{lstlisting}[style=cpp,title={\file{src/main.cpp}}]
#include "command.h"
#include "store.h"
#include <iostream>
#include <memory>
#include <string>

static void printUsage() {
    std::cerr << "Usage: kvstore <command> [args...]" << std::endl;
    std::cerr << "Commands:" << std::endl;
    std::cerr << "  set <key> <value>   Set a key-value pair"
              << std::endl;
    std::cerr << "  get <key>           Get value by key"
              << std::endl;
    std::cerr << "  delete <key>        Delete a key"
              << std::endl;
    std::cerr << "  list                List all pairs"
              << std::endl;
    std::cerr << "  count               Count entries"
              << std::endl;
    std::cerr << "  export <csv|json>   Export data"
              << std::endl;
}

static std::unique_ptr<Command> parseCommand(
        int argc, char* argv[]) {
    if (argc < 2) {
        return nullptr;
    }

    const std::string cmd = argv[1];

    if (cmd == "set") {
        if (argc != 4) {
            std::cerr << "Error: set requires <key> <value>"
                      << std::endl;
            return nullptr;
        }
        const std::string key = argv[2];
        if (key.empty()) {
            std::cerr << "Error: key cannot be empty"
                      << std::endl;
            return nullptr;
        }
        return std::make_unique<SetCommand>(key, argv[3]);
    }

    if (cmd == "get") {
        if (argc != 3) {
            std::cerr << "Error: get requires <key>"
                      << std::endl;
            return nullptr;
        }
        return std::make_unique<GetCommand>(argv[2]);
    }

    if (cmd == "delete") {
        if (argc != 3) {
            std::cerr << "Error: delete requires <key>"
                      << std::endl;
            return nullptr;
        }
        return std::make_unique<DeleteCommand>(argv[2]);
    }

    if (cmd == "list") {
        if (argc != 2) {
            std::cerr << "Error: list takes no arguments"
                      << std::endl;
            return nullptr;
        }
        return std::make_unique<ListCommand>();
    }

    if (cmd == "count") {
        if (argc != 2) {
            std::cerr << "Error: count takes no arguments"
                      << std::endl;
            return nullptr;
        }
        return std::make_unique<CountCommand>();
    }

    if (cmd == "export") {
        if (argc != 3) {
            std::cerr << "Error: export requires <csv|json>"
                      << std::endl;
            return nullptr;
        }
        auto exporter = createExporter(argv[2]);
        if (!exporter) {
            std::cerr << "Error: unknown format: " << argv[2]
                      << " (use csv or json)" << std::endl;
            return nullptr;
        }
        return std::make_unique<ExportCommand>(
            std::move(exporter));
    }

    std::cerr << "Error: unknown command: " << cmd
              << std::endl;
    return nullptr;
}

int main(int argc, char* argv[]) {
    auto cmd = parseCommand(argc, argv);
    if (!cmd) {
        printUsage();
        return 1;
    }

    KVStore store("kvstore.dat");
    store.load();
    cmd->execute(store);
    return 0;
}
\end{lstlisting}

% ============================================================
\chapter{패턴 요약표 (Pattern Quick Reference)}
\label{app:patterns}
% ============================================================

\begin{center}
\begin{longtable}{>{\bfseries}p{3cm} p{4cm} p{5.5cm}}
\toprule
패턴 & 핵심 개념 & 프로젝트 적용 \\
\midrule
\endhead

RAII &
생성자에서 자원 획득,\newline
소멸자에서 자원 해제 &
\code{KVStore} 클래스,\newline
\code{ifstream}/\code{ofstream} 파일 I/O \\
\midrule

값 의미론 &
복사가 기본, 참조는 명시적\newline
\code{const T\&}로 불필요한 복사 방지 &
모든 함수 매개변수에 적용 \\
\midrule

추상 클래스\newline(인터페이스) &
\code{virtual ... = 0} 순수 가상 함수\newline
다형성의 기반 &
\code{Command} 기반 클래스,\newline
\code{Exporter} 기반 클래스 \\
\midrule

Command &
요청을 객체로 캡슐화\newline
인자를 멤버로 저장 &
\code{SetCommand}, \code{GetCommand},\newline
\code{DeleteCommand} 등 \\
\midrule

Factory &
객체 생성 로직을 분리\newline
호출자가 구체 타입을 몰라도 됨 &
\code{parseCommand()} — CLI → Command,\newline
\code{createExporter()} — format → Exporter \\
\midrule

Strategy &
알고리즘을 인터페이스 뒤에 숨기고\newline
실행 시점에 교체 &
\code{CsvExporter}, \code{JsonExporter} \\
\midrule

스마트 포인터 &
\code{unique\_ptr}: 단일 소유권\newline
\code{std::move}: 소유권 이전 &
Factory 함수의 반환 타입,\newline
\code{ExportCommand}의 멤버 \\

\bottomrule
\end{longtable}
\end{center}

% ============================================================
\chapter{흔한 함정 (Common Pitfalls)}
\label{app:pitfalls}
% ============================================================

\section{댕글링 참조 (Dangling Reference)}

\begin{lstlisting}[style=cpp]
// 위험! 지역 변수의 참조를 반환
const std::string& bad() {
    std::string s = "hello";
    return s;  // s는 함수 끝에서 소멸 -> 댕글링 참조!
}

// 올바른 방법: 값으로 반환
std::string good() {
    std::string s = "hello";
    return s;  // 이동(move)이 자동으로 일어남
}
\end{lstlisting}

\begin{warningbox}
지역 변수의 참조(\code{\&})를 반환하지 마세요. 함수가 끝나면 지역 변수는 사라지고, 참조는 이미 없는 메모리를 가리킵니다. 값으로 반환하면 C++17 이상에서 이동 최적화(RVO/NRVO)가 자동으로 적용됩니다.
\end{warningbox}

\section{객체 슬라이싱 (Object Slicing)}

\begin{lstlisting}[style=cpp]
class Base {
public:
    virtual void hello() { std::cout << "Base\n"; }
};

class Derived : public Base {
public:
    void hello() override { std::cout << "Derived\n"; }
    int extra_ = 42;
};

Derived d;
Base b = d;  // 슬라이싱! extra_가 잘려나감
b.hello();   // "Base" 출력 (Derived가 아님!)

Base& ref = d;   // 참조: OK
ref.hello();     // "Derived" 출력

Base* ptr = &d;  // 포인터: OK
ptr->hello();    // "Derived" 출력
\end{lstlisting}

\begin{warningbox}
파생 클래스 객체를 기반 클래스 \textbf{값}으로 복사하면, 파생 부분이 잘려나갑니다. 다형성을 사용하려면 반드시 \textbf{참조}(\code{\&})나 \textbf{포인터}(\code{*}, \code{unique\_ptr})를 사용하세요.
\end{warningbox}

\section{암묵적 변환 (Implicit Conversion)}

\begin{lstlisting}[style=cpp]
class FilePath {
public:
    FilePath(const std::string& path) : path_(path) {}
    // explicit이 없으면...
private:
    std::string path_;
};

void process(FilePath fp) { /* ... */ }

process("hello.txt");  // string -> FilePath 자동 변환됨!
// 의도하지 않은 변환이 버그의 원인이 될 수 있음
\end{lstlisting}

\textbf{해결}: 인자 1개인 생성자에는 항상 \code{explicit}을 붙이세요.

\section{초기화 순서 (Initialization Order)}

멤버 변수는 \textbf{선언 순서}대로 초기화됩니다. 초기화 리스트의 순서가 아닙니다!

\begin{lstlisting}[style=cpp]
class Example {
    int a_;    // 1번째로 초기화
    int b_;    // 2번째로 초기화
public:
    // 경고: b_가 a_보다 먼저 초기화 리스트에 있지만,
    // 실제로는 a_가 먼저 초기화됨 (선언 순서)
    Example(int x) : b_(x), a_(b_) {}  // 위험! a_는 미초기화된 b_를 사용
};
\end{lstlisting}

\begin{warningbox}
\textbf{규칙}: 초기화 리스트의 순서를 선언 순서와 일치시키세요. 대부분의 컴파일러가 \code{-Wreorder} 경고를 줍니다.
\end{warningbox}

\section{헤더에서 using namespace std}

\begin{lstlisting}[style=cpp]
// store.h -- 절대 하지 마세요!
#ifndef STORE_H
#define STORE_H

using namespace std;  // 이 헤더를 include하는
                      // 모든 파일에 영향!

class KVStore {
    map<string, string> data_;  // std:: 생략 가능하지만...
};
#endif
\end{lstlisting}

\begin{warningbox}
헤더 파일에 \code{using namespace std;}를 쓰면 이 헤더를 포함하는 모든 파일의 이름 공간이 오염됩니다. 항상 \code{std::}를 명시적으로 쓰세요.
\end{warningbox}

% ============================================================

\end{document}
